When perusing for possible airfoils to base one's implementation of the unsteady vortex lattice method, one will often find simple \texttt{dat}-files, whose contents are the vertices of a polygon.
This polygon is a piecewise linear approximation of some airfoil curve, though the data is not always readily available for any and all airfoils at one's disposal.
For the insteady vortex lattice method, we want the mean camber line for the aerodynamic calculations and discretization into bound panels on the wind turbine blade.
Extracting the skeleton from the airfoil may be as simple as finding it analytically through splines.
It would be interesting, however, if one could find the skeleton using a more rudimentary method, like the one proposed in conjecture \ref{con:medial-axis}.
For future studies, a script for calculating the medial axis through the \textsc{Voronoj} diagram was implemented, yielding figure \ref{fig:karman-trefftz_voronoj} for the \textsc{Kármán}--\textsc{Trefftz} airfoil of figure \ref{subfig:karman-trefftz_foil}.

\begin{figure}[h]
  \centering
  \includegraphics[width = .75\linewidth]{karman-trefftz_voronoj.pdf}
  \caption{\textsc{Voronoj} diagram of the \textsc{von Kármán}--\textsc{Trefftz} airfoil (dashed).}
  \label{fig:karman-trefftz_voronoj}
\end{figure}

\noindent
We note that the \textsc{Voronoj} diagram is rather jagged towards the trailing edge, yielding in turn a jagged medial axis.
The medial axis is drawn with the analytically calculated skeleton of the airfoil in figure \ref{fig:karman-trefftz_medial_axis}

\begin{figure}[h]
  \centering
  \includegraphics[width = .75\linewidth]{karman-trefftz_medial_axis.pdf}
  \caption{Medial axis of the \textsc{von Kármán}--\textsc{Trefftz} airfoil (dashed), and the skeleton (full).}
  \label{fig:karman-trefftz_medial_axis}
\end{figure}

\noindent
With half the number of vertices, we attempt an even more naïve approach in figure \ref{fig:karman-trefftz_camber}.
From each vertex on the airfoil, we draw a vertical line to the opposite side of the polygon, and calculate the point midway between the vertex and the intersection on the opposite side.
This did indeed yield much better results.

\begin{figure}[h]
  \centering
  \includegraphics[width = .75\linewidth]{karman-trefftz_camber.pdf}
  \caption{Mean camber line (dashed) and skeleton (full) of the \textsc{von Kármán}--\textsc{Trefftz} airfoil.}
  \label{fig:karman-trefftz_camber}
\end{figure}

\noindent
For the \textsc{Selig}--\textsc{Donovan} airfoil SD7032, whose number of vertices is even lower than this, the \textsc{Voronoj} diagram was unfortunately not the better option between the two naïve methods.

\begin{figure}[h]
  \centering
  \includegraphics[width = .75\linewidth]{sd7032_voronoj.pdf}
  \caption{\textsc{Voronoj} diagram of SD7032 (dashed).}
  \label{fig:sd7032_voronoj}
\end{figure}


\begin{figure}[h]
  \centering
  \includegraphics[width = .75\linewidth]{sd7032_medial_axis.pdf}
  \caption{Medial axis of SD7032 (dashed) and mean camber line (full).}
  \label{fig:sd7032_medial_axis}
\end{figure}
